\paragraph{window-\texorpdfstring{$6$}{6} 的根--权交会刚性、orbit--Levi 骨架与二层 boundary 种子的 \texorpdfstring{$\ZZ_6$}{Z6} 主齐性}
在四步 uplift 同构、boundary-\(10\) 的根--Cartan 四块分裂、window-\(6\) 的 \(C_6\) 轨道骨架，以及二层 boundary 种子的两层覆盖已经固定之后，weight-layer、orbit--Levi、sheet parity 与等变编码四条离散接口可进一步压缩为下列四条结构性结论。其共同特征是：\(9+6+3+3\) 的 weight-layer 分块与 \(6+6+6+3\) 的根--Cartan 分块之间只有唯一的五原子共同细化；唯一三元轨道被锁定于 \(A_2^{(-)}\)；二层 boundary 种子形成正则 \(\ZZ_6\)-torsor；而任何保持全部内生对称的外置编码都必须真实含有阶 \(6\) 元。

\begin{theorem}[window-\texorpdfstring{$6$}{6} 的 weight-layer 分块与 \texorpdfstring{$B_3/C_3$}{B3/C3} 根分块的唯一五原子共同细化]\label{thm:xi-time-part53da-root-weight-common-refinement-five-atoms}\leanverified{paper\_xi\_time\_part53da\_root\_weight\_common\_refinement\_five\_atoms}
\[
X_6=X_6^{\mathrm{cyc}}\sqcup X_6^{\mathrm{bdry}},
\qquad
|X_6|=21,\quad |X_6^{\mathrm{cyc}}|=18,\quad |X_6^{\mathrm{bdry}}|=3,
\]
且
\[
X_6^{\mathrm{bdry}}
=
\{\texttt{100001},\ \texttt{100101},\ \texttt{101001}\}.
\]
记表 \ref{tab:window6_patisalam_9633_partition} 给出的 weight-layer 刚性分块为
\[
X_6^{\mathrm{cyc}}
=
X_6^{(2)}\sqcup X_6^{(1)}\sqcup X_6^{(L)},
\qquad
|X_6^{(2)}|=9,\quad |X_6^{(1)}|=6,\quad |X_6^{(L)}|=3,
\]
并以 \(X_6^{(R)}:=X_6^{\mathrm{bdry}}\) 补成
\[
X_6=X_6^{(2)}\sqcup X_6^{(1)}\sqcup X_6^{(L)}\sqcup X_6^{(R)}.
\]
另一方面，沿用定理 \ref{thm:boundary-shift4-uplift-isomorphism} 的 uplift 同构
\[
\iota_6:X_6\xrightarrow{\sim}X_{10}^{\mathrm{bdry}},
\qquad
\iota_6(v)=10\,v\,01,
\]
\[
X_6
=
X_6^{\mathrm{short}}
\sqcup
X_6^{A_2(+)}
\sqcup
X_6^{A_2(-)}
\sqcup
X_6^{\mathrm{Cartan}},
\]
其中 \(X_6^{\mathrm{short}},X_6^{A_2(+)},X_6^{A_2(-)},X_6^{\mathrm{Cartan}}\) 分别是定理 \ref{thm:boundary-m10-b3c3-rigid-four-block-split} 的四块分裂在 \(\iota_6\) 下的原像，故
\[
21=6+6+6+3.
\]
则这两种分块的共同细化恰有且仅有五个非空原子，其交会矩阵唯一为
\[
\begin{array}{c|cccc}
& \mathrm{short} & A_2^{(+)} & A_2^{(-)} & \mathrm{Cartan}\\
\hline
X_6^{(2)} & 0 & 3 & 6 & 0\\
X_6^{(1)} & 6 & 0 & 0 & 0\\
X_6^{(L)} & 0 & 3 & 0 & 0\\
X_6^{(R)} & 0 & 0 & 0 & 3
\end{array}
\]
从而共同细化的五个原子大小严格为
\[
6\ \oplus\ 3\ \oplus\ 6\ \oplus\ 3\ \oplus\ 3.
\]
\end{theorem}
\begin{proof}
表 \ref{tab:window6_patisalam_9633_partition} 与表 \ref{tab:boundary_uplift_m10_dictionary} 允许直接逐词核对。

首先，六个 weight-\(1\) 词
\[
\texttt{100000},\ \texttt{010000},\ \texttt{001000},\ \texttt{000100},\ \texttt{000010},\ \texttt{000001}
\]
恰对应六条 short roots，因此
\[
X_6^{(1)}=X_6^{\mathrm{short}}.
\]

其次，
\[
X_6^{(L)}=\{\texttt{000000},\ \texttt{010101},\ \texttt{101010}\}
\]
在根字典中全部落入 \(A_2^{(+)}\)，于是
\[
X_6^{(L)}\subset X_6^{A_2(+)}
\quad\text{且}\quad
|X_6^{(L)}\cap X_6^{A_2(+)}|=3.
\]

再者，九个 weight-\(2\) 词中
\[
\texttt{010100},\ \texttt{100010},\ \texttt{010001}
\]
落入 \(A_2^{(+)}\)，其余六个
\[
\texttt{101000},\ \texttt{100100},\ \texttt{010010},\ \texttt{001010},\ \texttt{001001},\ \texttt{000101}
\]
落入 \(A_2^{(-)}\)。因此
\[
|X_6^{(2)}\cap X_6^{A_2(+)}|=3,
\qquad
|X_6^{(2)}\cap X_6^{A_2(-)}|=6.
\]

最后，boundary 三词
\[
\texttt{100001},\ \texttt{100101},\ \texttt{101001}
\]
在表 \ref{tab:boundary_uplift_m10_dictionary} 中全部对应 Cartan 方向，因此
\[
X_6^{(R)}=X_6^{\mathrm{Cartan}}
\quad\text{且}\quad
|X_6^{(R)}\cap X_6^{\mathrm{Cartan}}|=3.
\]

四列总数已经分别达到 \(6,6,6,3\)，故除此五个非空交块外不再可能有新的非空交会。于是交会矩阵被唯一锁定，共同细化恰由五个非空原子组成。证毕。
\end{proof}

\begin{theorem}[orbit--Levi 骨架与根--Cartan 分裂的交会矩阵完全唯一]\label{thm:xi-time-part53da-orbit-levi-root-cartan-intersection}
沿用表 \ref{tab:window6_c6_orbit_decomposition} 的 \(C_6\) 轨道分解，并记
\[
X_6^{\mathrm{cyc},3}:=\mathcal{O}(\texttt{001001})
=
\{\texttt{001001},\texttt{010010},\texttt{100100}\},
\qquad
X_6^{\mathrm{cyc},15}:=X_6^{\mathrm{cyc}}\setminus X_6^{\mathrm{cyc},3}.
\]
则在
\[
X_6=X_6^{\mathrm{cyc},15}\sqcup X_6^{\mathrm{cyc},3}\sqcup X_6^{\mathrm{bdry}}
\qquad (21=15\oplus 3\oplus 3)
\]
与定理 \ref{thm:xi-time-part53da-root-weight-common-refinement-five-atoms} 的根--Cartan 分块之间，交会矩阵唯一为
\[
\begin{array}{c|cccc}
& \mathrm{short} & A_2^{(+)} & A_2^{(-)} & \mathrm{Cartan}\\
\hline
X_6^{\mathrm{cyc},15} & 6 & 6 & 3 & 0\\
X_6^{\mathrm{cyc},3}  & 0 & 0 & 3 & 0\\
X_6^{\mathrm{bdry}}   & 0 & 0 & 0 & 3
\end{array}
\]
并且同时在 orbit--Levi 骨架与根--Cartan 分块上常值的函数空间维数恰为 \(4\)。
\end{theorem}
\begin{proof}
表 \ref{tab:window6_c6_orbit_decomposition} 表明，唯一的三元轨道正是
\[
X_6^{\mathrm{cyc},3}
=
\{\texttt{001001},\texttt{010010},\texttt{100100}\}.
\]
再由表 \ref{tab:boundary_uplift_m10_dictionary}，这三个词分别对应
\[
(0,1,-1),\qquad (1,0,-1),\qquad (-1,1,0),
\]
因而全部落在 \(A_2^{(-)}\)。故
\[
X_6^{\mathrm{cyc},3}\subset X_6^{A_2(-)},
\qquad
|X_6^{\mathrm{cyc},3}\cap X_6^{A_2(-)}|=3.
\]

又 \(X_6^{\mathrm{bdry}}\) 的三个词在同一字典中全部对应 \(\mathrm{Cartan}\) 方向，因此
\[
|X_6^{\mathrm{bdry}}\cap X_6^{\mathrm{Cartan}}|=3.
\]

由定理 \ref{thm:xi-time-part53da-root-weight-common-refinement-five-atoms}，整个 cyclic 扇区与根--Cartan 字典的交会总数为
\[
6\ \mathrm{short}\ +\ 6\ A_2^{(+)}\ +\ 6\ A_2^{(-)}.
\]
从中扣除 \(X_6^{\mathrm{cyc},3}\) 已经占去的 \(3\) 个 \(A_2^{(-)}\) 方向，剩余的 \(X_6^{\mathrm{cyc},15}\) 便必然恰好分布为
\[
6\ \mathrm{short}\ +\ 6\ A_2^{(+)}\ +\ 3\ A_2^{(-)}.
\]
于是第一行也被唯一锁定，得到所示交会矩阵。

最后，一个函数同时在 orbit--Levi 三块与根--Cartan 四块上常值，当且仅当它在全部非空交块上常值。由于这里只有四个非空交块，这样的函数空间便恰为四维。特别地，唯一三元轨道被根字典强制置于 \(A_2^{(-)}\)，而不能落入 \(A_2^{(+)}\)。证毕。
\end{proof}

\begin{theorem}[window-\texorpdfstring{$6$}{6} 的二层 boundary 种子是正则 \texorpdfstring{$\ZZ_6$}{Z6}-torsor]\label{thm:xi-time-part53da-boundary-z6-torsor}
\leanverified{paper\_xi\_time\_part53da\_boundary\_z6\_torsor}
设
\[
X_6^{\mathrm{bdry}}
=
\{\texttt{100001},\ \texttt{100101},\ \texttt{101001}\},
\]
并沿用表 \ref{tab:fold6_boundary_sheet_pairs} 的 dyadic uplift 词典：每个 boundary 词在 \(\Fold_6^{\mathrm{bin}}\) 下恰有两个原像，且同一纤维内两点的差值恒为 \(34=F_9\)。记
\[
\widetilde X_6^{\mathrm{bdry}}
:=
\bigsqcup_{w\in X_6^{\mathrm{bdry}}}
\bigl(\Fold_6^{\mathrm{bin}}\bigr)^{-1}(w),
\qquad
|\widetilde X_6^{\mathrm{bdry}}|=6.
\]
则 \(\widetilde X_6^{\mathrm{bdry}}\) 上存在一个自由且传递的 \(\ZZ_6\)-作用。特别地，
\[
\CC[\widetilde X_6^{\mathrm{bdry}}]\cong \CC[\ZZ_6].
\]
作为复置换模，\(\CC[\widetilde X_6^{\mathrm{bdry}}]\) 恰为 \(\ZZ_6\) 的正则表示。
\end{theorem}
\begin{proof}
表 \ref{tab:fold6_boundary_sheet_pairs} 表明，\(\widetilde X_6^{\mathrm{bdry}}\) 由三条两点纤维组成，因此确有
\[
|\widetilde X_6^{\mathrm{bdry}}|=6.
\]
记
\[
\zeta:\widetilde X_6^{\mathrm{bdry}}\to \widetilde X_6^{\mathrm{bdry}}
\]
为每条两点纤维上的对角 sheet-flip。另一方面，三条 boundary 轴的循环置换给出一个 \(C_3\)-作用于基底三点集 \(X_6^{\mathrm{bdry}}\)，并自然提升为 \(\widetilde X_6^{\mathrm{bdry}}\) 上的置换作用。

由于 \(\zeta\) 只翻转 sheet 而不改变 boundary 轴，\(C_3\) 只循环 boundary 轴而不改变 sheet，二者可交换，故得到直积作用
\[
\langle\zeta\rangle\times C_3\curvearrowright \widetilde X_6^{\mathrm{bdry}}.
\]
再由
\[
\langle\zeta\rangle\times C_3
\cong
\ZZ_2\times \ZZ_3
\cong
\ZZ_6
\]
便可把该作用视为 \(\ZZ_6\)-作用。

该作用是传递的：给定任意两点 \(x,y\in \widetilde X_6^{\mathrm{bdry}}\)，先用 \(C_3\) 把 \(x\) 所在纤维送到 \(y\) 所在纤维；若此时 sheet 标号不同，再施加一次 \(\zeta\) 即可把像点送到 \(y\)。

另一方面，
\[
\bigl|\langle\zeta\rangle\times C_3\bigr|=2\cdot 3=6=|\widetilde X_6^{\mathrm{bdry}}|.
\]
有限群在同基数集合上的传递作用自动自由，故 \(\widetilde X_6^{\mathrm{bdry}}\) 是一个 \(\ZZ_6\)-torsor。最后，任一有限群在自身 torsor 上的置换表示都与其正则表示同构，因此
\[
\CC[\widetilde X_6^{\mathrm{bdry}}]\cong \CC[\ZZ_6].
\]
证毕。
\end{proof}

\begin{corollary}[保持全部内生对称的 boundary 编码必须含有阶六元]\label{cor:xi-time-part53da-equivariant-boundary-coding-order-six}
\leanverified{paper\_xi\_time\_part53da\_equivariant\_boundary\_coding\_order\_six}
沿用定理 \ref{thm:xi-time-part53da-boundary-z6-torsor} 的 \(\ZZ_6\)-作用。设 \(A\) 为有限阿贝尔群，则下述两件事等价：
\begin{enumerate}
  \item 存在单射 \(\Phi:\widetilde X_6^{\mathrm{bdry}}\to A\) 与群同态 \(\rho:\ZZ_6\to A\)，使得
  \[
  \Phi(g\cdot x)=\rho(g)+\Phi(x)
  \qquad
  (g\in \ZZ_6,\ x\in \widetilde X_6^{\mathrm{bdry}});
  \]
  \item \(A\) 含有一个循环子群同构于 \(\ZZ_6\)，等价地，\(A\) 含有一个阶 \(6\) 元。
\end{enumerate}
特别地：
\begin{enumerate}
  \item 对任意 \(r\ge 1\)，纯二进寄存器 \(A\cong (\ZZ/2\ZZ)^r\) 不可能承载这样的等变编码；
  \item 对任意 \(r\ge 1\)，纯三进寄存器 \(A\cong (\ZZ/3\ZZ)^r\) 也不可能承载这样的等变编码；
  \item 在全部有限阿贝尔目标中，最小可行者恰为 \(\ZZ_6\cong \ZZ_2\times \ZZ_3\)。
\end{enumerate}
\end{corollary}
\begin{proof}
把定理 \ref{thm:xi-time-part53da-boundary-z6-torsor} 中的作用群识别为 \(\ZZ_6\)。若存在这样的等变单射 \(\Phi\)，则必存在群同态
\[
\rho:\ZZ_6\to A
\]
使得
\[
\Phi(g\cdot x)=\rho(g)+\Phi(x)
\qquad
(g\in \ZZ_6,\ x\in \widetilde X_6^{\mathrm{bdry}}).
\]
固定任意 \(x_0\in \widetilde X_6^{\mathrm{bdry}}\)。由于 \(\widetilde X_6^{\mathrm{bdry}}\) 是 \(\ZZ_6\)-torsor，其轨道恰有 \(6\) 个点，故单射 \(\Phi\) 迫使
\[
\{\Phi(g\cdot x_0):\ g\in \ZZ_6\}
\]
也有 \(6\) 个不同元素。由等变性，
\[
\Phi(g\cdot x_0)=\rho(g)+\Phi(x_0),
\]
故该像轨道正是子群 \(\rho(\ZZ_6)\) 的一个平移副本；其大小只能是 \(|\rho(\ZZ_6)|\)。于是
\[
|\rho(\ZZ_6)|=6,
\]
从而 \(\rho(1)\in A\) 的阶恰为 \(6\)。

反之，若 \(A\) 含有一个阶 \(6\) 元 \(a\)，则取同态
\[
\rho:\ZZ_6\to A,\qquad \rho(1)=a.
\]
固定基点 \(x_0\in \widetilde X_6^{\mathrm{bdry}}\)，并对任意 \(g\in \ZZ_6\) 定义
\[
\Phi(g\cdot x_0):=\rho(g).
\]
由于 \(\widetilde X_6^{\mathrm{bdry}}\) 上的 \(\ZZ_6\)-作用自由且传递，该定义良好并且自动等变；又因 \(a\) 的阶为 \(6\)，\(\rho\) 为单射，故 \(\Phi\) 亦为单射。

最后，\((\ZZ/2\ZZ)^r\) 的每个元素阶都整除 \(2\)，\((\ZZ/3\ZZ)^r\) 的每个元素阶都整除 \(3\)，因此都不含阶 \(6\) 元；而 \(\ZZ_6\cong\ZZ_2\times\ZZ_3\) 显然含有阶 \(6\) 元，且阶数最小。证毕。
\end{proof}

\noindent
由此，window-\(6\) 的 weight-layer 分块、orbit--Levi 骨架、根--Cartan 字典与二层 boundary 覆盖并非彼此松散兼容的四组离散证书，而是由唯一交会矩阵与正则 \(\ZZ_6\)-torsor 共同锁定的一组刚性接口。于是，定理 \ref{thm:xi-boundary-uplift-mixed-radix-register-lower-bound} 给出的 mixed-radix 状态数下界，在这一最小窗口上可进一步强化为严格的群论判据：任何保留全部 boundary 内生对称的目标寄存器都必须真实含有阶 \(6\) 元。

\endinput
